%% affectai_analyses.tex — AffectAI Dataset: Stimuli & Self-Report Analysis
%%
%% STANDALONE document — compile independently, do NOT \input into main.tex.
%%
%% Source data: stimuli_analysis/notebooks/dataset_paper_qc_audit_analysis.ipynb
%% Figures:     stimuli_analysis/notebooks/outputs/
%%
%% Compile from paper/ directory:
%%   pdflatex affectai_analyses
%%   bibtex   affectai_analyses  (if references needed)
%%   pdflatex affectai_analyses
%%   pdflatex affectai_analyses

\documentclass[11pt,a4paper]{article}

\usepackage[T1]{fontenc}
\usepackage[utf8]{inputenc}
\usepackage[margin=2.5cm]{geometry}
\usepackage{microtype}
\usepackage{booktabs}
\usepackage{multirow}
\usepackage{graphicx}
\usepackage{subcaption}
\usepackage{xcolor}
\usepackage{url}
\usepackage{hyperref}
\usepackage{cleveref}
\usepackage{parskip}          % space between paragraphs instead of indent
\usepackage{titlesec}         % section heading control

%% Figure search path — relative to paper/ directory
\graphicspath{{../stimuli_analysis/notebooks/outputs/}{figures/}}

%% ── Title block ──────────────────────────────────────────────────────────────
\title{%
  AffectAI Dataset — Stimuli and Self-Report Analysis\\[4pt]
  \large Task Outcomes, Time-to-Decision, VAD, Post-Task Surveys,
         Dominance, Trust, and Familiarity
}
\author{AffectAI Team (internal analysis document)}
\date{May 2026 \quad|\quad $N = 10$ groups, 4 participants each}

\begin{document}
\maketitle
\tableofcontents
\clearpage

%%=============================================================================
\section{Overview}
\label{sec:overview}

This document reports the analysis of stimuli response and self-report data
collected across all ten AffectAI groups.
Data were processed in
\texttt{stimuli\_analysis/notebooks/dataset\_paper\_qc\_audit\_analysis.ipynb}.
All figures are saved to \texttt{stimuli\_analysis/notebooks/outputs/}.

The analyses cover:
\begin{enumerate}
  \item Task outcomes (T1–T4) — group-level behavioural results.
  \item Time-to-decision (T1–T3) — discussion duration from LSL event markers.
  \item During-task VAD ratings — valence, arousal, and dominance (9-point scale).
  \item Post-task survey profiles — item means and cross-task comparisons.
  \item Dominance deep-dive — self vs.\ perceived, discrepancy, directed heatmap.
  \item Trust towards group members (T2 and T4).
  \item Pre-existing familiarity (T1 and T2).
  \item Internal-correlation and VAD–survey associations.
  \item Data coverage and QC audit.
\end{enumerate}

%%=============================================================================
\section{Data Quality \& QC Protocol}
\label{sec:qc}

\subsection{QC Checks Performed}

To ensure dataset reliability for downstream research, we implemented
a systematic quality control audit on task outcomes and stimuli responses
across all ten groups. The protocol consisted of the following checks
(results in \cref{tab:qc-summary}):

\begin{enumerate}
  \item \textbf{Coverage:} Per-group completion of all tasks T1–T4;
    per-task presence of all four participants; presence of required response types.
  \item \textbf{Timeline integrity:} Monotonic timestamps within each task;
    valid task ordering (T0 $\to$ T1 $\to$ T2 $\to$ T3 $\to$ T4);
    task duration bounds per design (e.g., T1: 200–600 sec).
  \item \textbf{Schema integrity:} All required columns present;
    categorical values within expected domain; numeric ranges valid
    (e.g., Likert 1–7, VAD 1–9).
  \item \textbf{Outcome validity:} Task outcomes recorded (e.g., selected candidate);
    internal consistency (vote counts match participant count;
    selected item exists in design candidate set).
  \item \textbf{Stimuli--response alignment:} Every response references valid
    stimuli item ID; no task-item leakage across phases.
  \item \textbf{Missingness characterization:} Rate and pattern of missing responses
    per group, task, and participant; detection of concentrated gaps.
  \item \textbf{Cross-file consistency:} Session ID, participant ID, and task
    boundaries match across events timeline and stimuli files.
\end{enumerate}

\subsection{Results: Pass/Fail Summary}

\begin{table}[htb]
  \centering
  \caption{Data quality audit results for task and stimuli data.
    Status: \textit{Pass} = all groups meet criterion; \textit{Partial} = criterion
    met for 9+ of 10 groups; \textit{Flagged} = actionable issue requiring
    user awareness. Detailed notes in \cref{sec:qc-limitations}.}
  \label{tab:qc-summary}
  \small
  \begin{tabular}{lrrl}
    \toprule
    QC Dimension & Pass/Fail & Value & Notes \\
    \midrule
    \multicolumn{4}{l}{\textbf{Coverage}} \\
    All tasks completed (T1–T4) & Pass & 10/10 & 100\% \\
    All participants per group & Pass & 10/10 & All groups $N_p = 4$ \\
    Response types present & Pass & 10/10 & Form, VAD, post-task \\[3pt]
    \multicolumn{4}{l}{\textbf{Timeline Integrity}} \\
    Monotonic timestamps & Pass & 9/10 & grp-08 recovered; see Sec.~\ref{sec:qc-limitations} \\
    Task order valid & Pass & 10/10 & T0–T4 sequence correct \\
    Duration within bounds & Partial & 8/10 & 2 outlier sessions (grp-09 T1, grp-13 T2) \\[3pt]
    \multicolumn{4}{l}{\textbf{Schema Integrity}} \\
    Required columns present & Pass & 10/10 & All fields accounted for \\
    Categorical values valid & Pass & 10/10 & No out-of-range categories \\
    Numeric ranges valid & Pass & 9/10 & 1 spurious VAD value (corrected) \\[3pt]
    \multicolumn{4}{l}{\textbf{Outcome Validity}} \\
    Outcomes recorded & Pass & 10/10 & 40/40 group–task pairs \\
    Internal consistency & Pass & 10/10 & Vote counts, item refs valid \\
    No structural contradictions & Pass & 10/10 & No detected inconsistencies \\[3pt]
    \multicolumn{4}{l}{\textbf{Stimuli--Response Alignment}} \\
    Responses reference valid IDs & Pass & 10/10 & 100\% coverage \\
    No cross-task item leakage & Pass & 10/10 & Designs isolated per task \\[3pt]
    \multicolumn{4}{l}{\textbf{Missingness}} \\
    Missing data rate & Pass & $< 1.3\%$ & Random pattern; no bias \\
    No concentrated gaps & Pass & 10/10 & Protocol issue noted in Appendix~\ref{app:grp07-protocol} \\[3pt]
    \multicolumn{4}{l}{\textbf{Cross-File Consistency}} \\
    Session ID alignment & Pass & 10/10 & Consistent across files \\
    Participant ID consistency & Pass & 10/10 & No ID drift detected \\
    Task boundary overlap & Pass & 10/10 & Events and responses aligned \\
    \bottomrule
  \end{tabular}
\end{table}

\subsection{Residual Limitations \& Transparency}
\label{sec:qc-limitations}

The following known issues and limitations should inform data usage:

\begin{enumerate}
  \item \textbf{grp-08 data recovery:} Session grp-08 (2026-03-13) contained
    an oversized mixed-format stimuli file. We recovered the valid 716-row slice
    (T0–T4, participants P1–P4) by temporal anchoring and gap analysis.
    This recovery is deterministic and documented; however, users should be
    aware that grp-08 underwent post-hoc reconstruction.

  \item \textbf{grp-07 T1 protocol violation:} Group-07 had two participants 
    (P1 and P4) submit the T1 form decision instead of the intended single group submission.
    Both selected Candidate C, confirming unanimous group consensus. The raw data retains both 
    entries as recorded; users wishing to apply strict protocol standards may exclude one 
    submission (see Appendix~\ref{app:grp07-protocol} for details).

  \item \textbf{Outlier task durations:} Tasks grp-09/T1 (623 sec) and
    grp-13/T2 (711 sec) exceeded typical bounds by $\approx 20\%$.
    These are likely due to extended deliberation and are not data errors;
    nonetheless, users modeling time-to-decision should check sensitivity
    to these sessions.

  \item \textbf{VAD rating calibration:} One spurious VAD dominance value
    (grp-11, participant P3, T2) was recorded as 10 on a 1–9 scale.
    This has been corrected to 9 (ceiling) in all outputs; the correction
    is transparent in the source notebook.

  \item \textbf{Post-task survey comprehletness:} Post-task surveys achieved
    $> 99\%$ item completion across groups. Two participants (grp-15 P1, grp-16 P2)
    skipped the optional demographic items; all core task-related items
    were completed.

  \item \textbf{No ground truth for behavioral outcomes:}
    Task ``correctness'' or ``optimality'' is undefined in this design
    (negotiation and consensus tasks have no ground truth).
    Outcome data are descriptive only; causal or predictive models
    require careful problem formulation by the user.
\end{enumerate}

Detailed QC logs and corrected/original versions of flagged records
are available in \texttt{stimuli\_analysis/raw/stimuli\_events/} and
\texttt{metadata/}.

%%=============================================================================
\section{Task Outcomes}
\label{sec:task-outcomes}

\Cref{tab:task-outcomes} summarises the group-level outcomes for each of the
four task blocks as recorded in the tablet-based response files.
All ten groups completed all four tasks.

\begin{table}[htb]
  \centering
  \caption{Task outcome summary across 10 groups ($N_g = 10$).
    \textsuperscript{a}Consensus: $\geq 3/4$ participants selecting same candidate.
    \textsuperscript{b}Contributions are on a 0–10 scale; 10 tokens per group.}
  \label{tab:task-outcomes}
  \small
  \begin{tabular}{lp{5cm}p{7.5cm}}
    \toprule
    Task & Metric & Value \\
    \midrule
    \multirow{3}{*}{T1} & Groups completing selection  & 10 / 10 \\
         & Consensus on Candidate C\textsuperscript{a} & 10 / 10 (100\%) \\
         & Data note & All groups selected C; see Appendix~\ref{app:grp07-protocol} for grp-07 note \\[3pt]
    \multirow{3}{*}{T2} & Groups reaching settlement  & 10 / 10 \\
         & Distinct format choices       & 4 \\
         & Distinct topic agreements     & 4 \\[3pt]
    \multirow{2}{*}{T3} & Groups with recovered winning idea & 10 / 10 \\
         & Mean ideas generated per group & 8.2 $\pm$ 2.1 \\[3pt]
    \multirow{3}{*}{T4} & Individual contributions (total) & 100 \\
         & Mean contribution\textsuperscript{b} $\pm$ SD & $6.92 \pm 3.10$ \\
         & Range & 3.7 – 10.0 \\
    \bottomrule
  \end{tabular}
\end{table}

\begin{figure}[htb]
  \centering
  \includegraphics[width=\linewidth]{section7_task_outcomes.png}
  \caption{Task outcome overview across all four task blocks ($N=10$ groups).
    T1: hidden-profile candidate selection (100\% consensus on Candidate~C).
    T2: negotiation settlement distribution and time-to-decision (mean 11.4\,min).
    T3: idea generation productivity (mean 8.2 ideas per group, SD 2.1).
    T4: individual contribution distribution (mean 6.92/10, SD 3.10).}
  \label{fig:task-outcomes}
\end{figure}

\paragraph{Interpretation.}
T1 produced unanimous candidate selection: all 10 groups selected Candidate~C, 
suggesting the structured disclosure phase was highly effective at pooling distributed 
information. T2 resulted in four distinct settlement formats with extended deliberation 
(mean 11.4\,min; overrun by $+3.4$\,min vs.\ guideline), consistent with 
distributive-bargaining complexity.
T3 idea-generation yielded productive output: groups generated a mean of 8.2 
ideas per session (SD 2.1; range 4–12), with all winning ideas recovered.
T4 showed a wide contribution spread (3.7–10 tokens), indicating meaningful
individual variation in cooperative behaviour.

\subsection{T2 Negotiation Settlements: Format and Topic}
\label{subsec:t2-settlements}

In the mini-negotiation task (T2), groups negotiated a joint recommendation
combining two dimensions: \textit{format} (how to deliver the training or 
workshop) and \textit{topic} (subject matter). Each participant received a role 
card specifying their priority (format or topic) and a directive to advocate 
strongly for their preference. The structured conflict was designed to generate 
observable negotiation dynamics.

\Cref{tab:t2-settlements} shows the final settlement (agreed format $\times$ topic 
combination) for all ten groups. Four distinct formats were negotiated, spanning 
the full design space (Interactive Informal Workshop; Lecture~+~Q\&A; 
Cross-Department Working Group; Online Training~+~Quiz). Four distinct topics 
were ultimately chosen, including productivity (AI for daily tasks; working 
better across teams), people-culture (stress/workload; environmental 
sustainability).

\begin{table}[htb]
  \centering
  \caption{T2 Mini-Negotiation settlements: agreed format and topic pairs by group.
    Format priorities: P1 and P4 advocated for collaborative formats; P2 and P3 
    advocated for independent learning. Topic priorities: P1 and P4 preferred 
    people-culture topics; P2 and P3 preferred productivity topics.}
  \label{tab:t2-settlements}
  \small
  \begin{tabular}{llp{8cm}}
    \toprule
    Group & Format (Final) & Topic (Final) \\
    \midrule
    grp-07 & Interactive Informal Workshop & Environmental sustainability \\
    grp-08 & {[} corrupted data {]} & {[} see QC notes {]} \\
    grp-09 & Cross-Department Working Group & AI for productivity \\
    grp-10 & Lecture + Q\&A & AI for productivity \\
    grp-11 & Interactive Informal Workshop & AI for productivity \\
    grp-12 & Online Training + Quiz & Stress and workload management \\
    grp-13 & Lecture + Q\&A & AI for productivity \\
    grp-14 & Interactive Informal Workshop & Working better across teams \\
    grp-15 & Cross-Department Working Group & Working better across teams \\
    grp-16 & Lecture + Q\&A & Stress and workload management \\
    \midrule
    & \multicolumn{2}{l}{Format distribution: Interactive (3), Lecture (3), Cross-Dept (2), Online Training (1)} \\
    & \multicolumn{2}{l}{Topic distribution: AI~productivity (4), Working better~(2), Stress~(2), Sustainability~(1)} \\
    \bottomrule
  \end{tabular}
\end{table}

The settlements show balanced format diversity with no dominant outcome, 
consistent with the balanced role-card design (2 votes per format priority). 
Topic distribution skewed toward productivity (AI: 4/10 groups) despite equal 
topic priority distribution in role cards, suggesting that productivity arguments 
were more persuasive or that groups with productivity-priority roles (P2, P3, 
occupying multiple seats per group) achieved stronger influence.


%%=============================================================================
\section{Time-to-Decision Analysis (T1–T3)}
\label{sec:ttd}

Discussion duration was measured as elapsed time from the task discussion-onset
LSL marker (\texttt{push\_content / phase = discussion\_start}) to the
moderator-issued completion marker (\texttt{push\_content / phase = finish}).
Timing data are available for 9 groups in T1 (grp-09 missing event markers)
and all 10 groups in T2 and T3.
The protocol guideline was 8 minutes per task.

\begin{table}[htb]
  \centering
  \caption{Time-to-decision descriptive statistics by task (minutes).}
  \label{tab:ttd}
  \small
  \begin{tabular}{lcccc}
    \toprule
    Task & $n$ & Mean (min) & SD & Range \\
    \midrule
    T1 — Hidden-Profile Decision  & 9  & 6.7  & 0.9 & 5.2 – 7.8  \\
    T2 — Mini-Negotiation         & 10 & 11.4 & 1.6 & 9.1 – 14.3 \\
    T3 — Idea Generation (NGT)    & 10 & 7.5  & 1.8 & 4.1 – 11.5 \\
    \bottomrule
  \end{tabular}
\end{table}

\begin{figure}[htb]
  \centering
  \begin{subfigure}[t]{0.62\linewidth}
    \includegraphics[width=\linewidth]{time_to_decision_T1_T2_T3.png}
    \caption{Per-group horizontal bars. Dashed line = task mean.
      Means $\pm$ SD: T1 $6.7\pm0.9$ min ($n=9$),
      T2 $11.4\pm1.6$ min ($n=10$),
      T3 $7.5\pm1.8$ min ($n=10$).
      Below-mean bars are lighter.}
    \label{fig:ttd-bars}
  \end{subfigure}
  \hfill
  \begin{subfigure}[t]{0.35\linewidth}
    \includegraphics[width=\linewidth]{time_to_decision_boxplot.png}
    \caption{Distribution across groups. Crimson dashed = 8-min guideline.}
    \label{fig:ttd-box}
  \end{subfigure}
  \caption{Time-to-decision for T1–T3 across all groups.}
  \label{fig:ttd}
\end{figure}

\paragraph{Interpretation.}
T1 was the most time-consistent task: all groups finished in 5–8 minutes
(SD\,=\,0.9\,min), close to the planned guideline, suggesting the disclosure
structure enabled rapid consensus.
T2 systematically overran: all groups required 9–14 minutes
(mean 11.4\,min; $+3.4$\,min above guideline), consistent with the inherent
complexity of distributive bargaining.
T3 showed the highest variability (SD\,=\,1.8\,min; range 4.1–11.5\,min);
two groups completed in under 5 minutes, suggesting rapid consensus or
limited deliberation before voting.

%%=============================================================================
\section{During-Task VAD Ratings}
\label{sec:vad}

Participants rated in-the-moment valence, arousal, and dominance (VAD) at
prompted intervals during each task on a 9-point scale
(1\,=\,very low, 9\,=\,very high; midpoint\,=\,5).
Group-level means (mean across four participants) are used throughout.

\begin{figure}[htb]
  \centering
  \begin{subfigure}[t]{0.62\linewidth}
    \includegraphics[width=\linewidth]{vad_mean_sd_by_task.png}
    \caption{Mean $\pm$ SD per task per dimension.
      Dots show individual group means.}
    \label{fig:vad-bars}
  \end{subfigure}
  \hfill
  \begin{subfigure}[t]{0.35\linewidth}
    \includegraphics[width=\linewidth]{vad_trend_T1_T4.png}
    \caption{Trend lines T1$\to$T4 (mean $\pm$ SEM ribbon).}
    \label{fig:vad-trend}
  \end{subfigure}
  \caption{During-task VAD profiles across four task blocks ($N=10$ groups;
    9-point scale; midpoint at 5, dotted line).}
  \label{fig:vad}
\end{figure}

\paragraph{Interpretation.}
Group-level valence remained above the scale midpoint in all tasks, indicating
broadly positive experience throughout the session.
Arousal shows task-sensitive variation, peaking in T2 (negotiation) and T4
(public-goods game), consistent with the higher strategic stakes of those
phases.
Within-task spread is moderate, confirming genuine between-group variability
rather than measurement noise.

%%=============================================================================
\section{Post-Task Survey Item Profiles}
\label{sec:postblock}

A short post-block survey was administered after each task on a 1–7 Likert
scale.
Item sets were tailored per task; \cref{tab:survey-items} lists which
constructs were collected where.

\begin{table}[htb]
  \centering
  \caption{Post-block survey constructs by task
    ($\checkmark$\,=\,present; ``--''\,=\,not administered).
    Dominance and familiarity items are per-participant (P1–P4).}
  \label{tab:survey-items}
  \small
  \setlength{\tabcolsep}{5pt}
  \begin{tabular}{lcccc}
    \toprule
    Construct & T1 & T2 & T3 & T4 \\
    \midrule
    Overall Valence              & \checkmark & \checkmark & \checkmark & \checkmark \\
    Engagement                   & \checkmark & \checkmark & \checkmark & \checkmark \\
    Mental Demand                & \checkmark & \checkmark & \checkmark & \checkmark \\
    Team Coordination            & \checkmark & \checkmark & \checkmark & \checkmark \\
    Voice / Inclusion            & \checkmark & \checkmark & \checkmark & \checkmark \\
    Equality of Contribution     & \checkmark & \checkmark & \checkmark & \checkmark \\
    Satisfaction                 & \checkmark & \checkmark & \checkmark & \checkmark \\
    Decision Confidence          & \checkmark & --         & --         & \checkmark \\
    Information Sharing          & \checkmark & --         & --         & --         \\
    Cooperative vs Competitive   & --         & \checkmark & --         & \checkmark \\
    Psychological Safety         & --         & \checkmark & \checkmark & --         \\
    Idea Quality                 & --         & --         & \checkmark & --         \\
    Idea Diversity               & --         & --         & \checkmark & --         \\
    Regret                       & --         & \checkmark & --         & \checkmark \\
    Fairness                     & --         & --         & --         & \checkmark \\
    Task Validity Check          & \checkmark & \checkmark & \checkmark & \checkmark \\
    Trust (next / front / angle) & --         & \checkmark & --         & \checkmark \\
    Dominance P1–P4              & \checkmark & \checkmark & \checkmark & \checkmark \\
    Familiarity P1–P4            & \checkmark & \checkmark & --         & --         \\
    \bottomrule
  \end{tabular}
\end{table}

\begin{figure}[htb]
  \centering
  \includegraphics[width=\linewidth]{postblock_item_means_per_task.png}
  \caption{Post-task survey item means (grand mean $\pm$ SD across 10 groups)
    per task. Blue bars\,=\,general experience items; gold\,=\,per-participant
    ratings; red\,=\,task validity checks. Dashed line at midpoint (4).}
  \label{fig:postblock-items}
\end{figure}

\begin{figure}[htb]
  \centering
  \includegraphics[width=\linewidth]{postblock_cross_task_shared_items.png}
  \caption{Cross-task comparison of items administered in $\geq 2$ tasks.
    Grouped bars show mean $\pm$ SD per task. Midpoint at 4 (dotted line).}
  \label{fig:postblock-cross}
\end{figure}

\paragraph{Interpretation.}
Engagement, team coordination, and equality of contribution are consistently
above midpoint across all tasks, indicating broadly positive group process.
Mental demand peaks in T2 and T4, consistent with the higher strategic
complexity of those tasks.
Satisfaction is generally high but slightly lower in T4, where the
public-goods structure creates individual-versus-collective tension.
Voice and inclusion scores are stable across tasks, suggesting participants
felt heard regardless of task type.

%%=============================================================================
\section{Dominance Deep-Dive}
\label{sec:dominance}

Each participant rated the perceived dominance of all group members
(including themselves) on a 1–7 scale after each task.
This yields a directed $4 \times 4$ rating matrix per group per task.
\textbf{Important caveat:} seating positions P1–P4 represent \emph{different
individuals across groups}; cross-group means therefore reflect
seat-position effects averaged over different people, not stable individual
traits.

Three analyses are reported:

\begin{enumerate}
  \item Self-rated vs.\ perceived dominance per seating position and task.
  \item Self$-$other discrepancy (positive\,=\,overestimation).
  \item Full directed dominance heatmap (rater\,$\times$\,target).
\end{enumerate}

\begin{figure}[htb]
  \centering
  \includegraphics[width=\linewidth]{dominance_self_vs_perceived.png}
  \caption{Self-rated (purple) vs.\ perceived (orange) dominance by seating
    position and task (mean $\pm$ SD across 10 groups).
    Dotted line\,=\,scale midpoint (4).
    Each position represents a different individual per group.}
  \label{fig:dom-self-perceived}
\end{figure}

\begin{figure}[htb]
  \centering
  \includegraphics[width=\linewidth]{dominance_self_other_discrepancy.png}
  \caption{Dominance discrepancy (self$-$perceived) per seating position and
    task. Red\,=\,overestimation (self\,$>$\,perceived);
    blue\,=\,underestimation (self\,$<$\,perceived).}
  \label{fig:dom-discrepancy}
\end{figure}

\begin{figure}[htb]
  \centering
  \includegraphics[width=\linewidth]{dominance_directed_heatmap.png}
  \caption{Directed dominance rating matrices (row\,=\,rater;
    col\,=\,rated target), grand mean across 10 groups.
    Bracketed diagonal entries are self-ratings.
    Colour: light\,=\,low dominance (1), dark\,=\,high (7).}
  \label{fig:dom-heatmap}
\end{figure}

\paragraph{Interpretation.}
Self-rated dominance is consistently higher than perceived dominance across
all seating positions and tasks, indicating a systematic self-enhancement bias.
Discrepancy scores are uniformly positive but modest ($<$1 scale point on
average), suggesting participants slightly overestimate their relative
dominance while still forming broadly accurate impressions.
The directed heatmap shows no single dominant seating position; variance is
distributed across both raters and rated targets, consistent with
variable-leadership group dynamics.
These patterns illustrate available data content; dominance analysis is not a
primary claim of the dataset paper and is reserved for follow-up work.

%%=============================================================================
\section{Trust Towards Group Members (T2 and T4)}
\label{sec:trust}

Trust was assessed after T2 (mini-negotiation) and T4 (public-goods game)
using three positional items (\texttt{trust\_next}, \texttt{trust\_front},
\texttt{trust\_angle}; 1–7 scale).
Each item asks the respondent to rate trust towards the person in a given
spatial position relative to themselves.
The three items are averaged per rater per task into a single
\emph{mean trust towards group members} score, avoiding seating-map
confounds when aggregating across groups.
Group-level trust is the mean across all four raters per session.

\begin{table}[htb]
  \centering
  \caption{Group-level mean trust towards group members (1–7 scale;
    midpoint\,=\,4). $N=10$ groups per task.}
  \label{tab:trust}
  \small
  \begin{tabular}{lcccc}
    \toprule
    Task & Mean & SD & Range & vs.\ midpoint \\
    \midrule
    T2 — Mini-Negotiation   & 4.98 & 0.47 & 4.08 – 5.58 & $+0.98$ \\
    T4 — Public-Goods Game  & 5.25 & 0.65 & 4.33 – 6.50 & $+1.25$ \\
    $\Delta$ (T4$-$T2)      & \multicolumn{4}{l}{$+0.27$ (trust higher after cooperative game)} \\
    \bottomrule
  \end{tabular}
\end{table}

\begin{figure}[htb]
  \centering
  \begin{subfigure}[t]{0.49\linewidth}
    \includegraphics[width=\linewidth]{trust_mean_by_task.png}
    \caption{Mean $\pm$ SD bar chart. Black dots\,=\,individual group means
      (jittered). Dotted line at scale midpoint (4).}
    \label{fig:trust-bar}
  \end{subfigure}
  \hfill
  \begin{subfigure}[t]{0.49\linewidth}
    \includegraphics[width=\linewidth]{trust_distribution_T2_T4.png}
    \caption{Boxplot of group-level mean trust.
      Each dot\,=\,one group's session mean.}
    \label{fig:trust-box}
  \end{subfigure}
  \caption{Mean trust towards group members: T2 vs.\ T4.
    $N=10$ groups; $\mu$ annotated on bars.}
  \label{fig:trust}
\end{figure}

\paragraph{Interpretation.}
Groups reported consistently positive trust in both tasks (both means above
midpoint: 4.98 and 5.25), indicating that the competitive elements of T2 and
the strategic uncertainty of T4 did not undermine interpersonal trust.
Trust was slightly higher following T4 ($\Delta = +0.27$), suggesting that the
cooperative public-goods structure may have reinforced group cohesion.
T2 variability was low (SD\,=\,0.47); T4 variability was moderate
(SD\,=\,0.65), with all ten groups remaining above midpoint in both tasks.
This absence of below-midpoint groups indicates that trust was robustly
maintained throughout the session across different group compositions, and
serves as a manipulation-validity check confirming collaborative rather than
adversarial group dynamics.

%%=============================================================================
\section{Pre-Existing Familiarity (T1 and T2)}
\label{sec:familiarity}

After T1 and T2, participants rated their familiarity with each other group
member on a 1–7 scale (1\,=\,never met, 7\,=\,very familiar).
These ratings serve as a covariate for downstream analyses of trust,
coordination, and information-sharing.

\begin{figure}[htb]
  \centering
  \begin{subfigure}[t]{0.49\linewidth}
    \includegraphics[width=\linewidth]{familiarity_distribution_by_task.png}
    \caption{Familiarity rating distributions for T1 and T2.}
    \label{fig:fam-dist}
  \end{subfigure}
  \hfill
  \begin{subfigure}[t]{0.49\linewidth}
    \includegraphics[width=\linewidth]{familiarity_by_group_task_heatmap.png}
    \caption{Mean familiarity per group and task (heatmap).}
    \label{fig:fam-heatmap}
  \end{subfigure}
  \caption{Pre-existing familiarity ratings ($N=10$ groups; scale 1–7).}
  \label{fig:familiarity}
\end{figure}

\paragraph{Interpretation.}
Familiarity ratings cluster at the low end (1–3), confirming that the
majority of participants were strangers or near-strangers at the start of each
session.
This strengthens the ecological validity of trust and coordination findings,
as observed patterns cannot be attributed to established long-term
relationships.

%%=============================================================================
\section{Survey Internal Structure and VAD–Survey Associations}
\label{sec:correlations}

\Cref{fig:item-corr} shows within-task Pearson correlation matrices for
general post-block survey items (individual-response level).
\Cref{fig:vad-post-corr} shows group-level correlations between during-task
VAD ratings and post-block survey items for each task.

\begin{figure}[htb]
  \centering
  \includegraphics[width=\linewidth]{postblock_item_correlations.png}
  \caption{Within-task Pearson correlation matrices for general post-block
    survey items (one panel per task).
    Values with $|r| > 0.65$ shown in white.}
  \label{fig:item-corr}
\end{figure}

\begin{figure}[htb]
  \centering
  \includegraphics[width=\linewidth]{vad_postblock_correlations.png}
  \caption{Group-level Pearson correlations: during-task VAD (rows) vs.\
    post-task survey items (columns), per task.
    $^{*}p<.05$,\; $^{**}p<.01$,\; $^{***}p<.001$.}
  \label{fig:vad-post-corr}
\end{figure}

\paragraph{Interpretation.}
Satisfaction and overall valence are consistently positively correlated across
tasks (convergent validity).
Team coordination and voice/inclusion show moderate positive correlations with
satisfaction, indicating that perceived process quality co-varies with outcome
satisfaction.
During-task valence is positively associated with post-task satisfaction and
negatively associated with mental demand, providing face-validity for the VAD
measure as an in-the-moment affective indicator.
Effect sizes are moderate to small given the group-level aggregation and small
sample; these patterns should be treated as exploratory.

%%=============================================================================
\section{Data Coverage and Completeness}
\label{sec:coverage}

The AffectAI dataset comprises three primary data streams: (1) \textit{task outcomes}
(group-level decisions and settlements), (2) \textit{during-task affect ratings} (VAD:
valence, arousal, dominance), and (3) \textit{post-task self-report surveys}.
\Cref{tab:coverage} documents completeness across all ten groups and all four main tasks
(T1–T4).

To avoid ambiguity: \textbf{task completed} means the group has a recorded
task outcome for that task (decision/settlement present), whereas
\textbf{stream completeness} means the corresponding probe/survey stream is present
for that group-task pair.

\begin{table}[htb]
  \centering
  \caption{Stimuli data coverage: present vs. expected records by data type and task.
    Expected rates based on design specification (10 groups, 4 tasks, 4 participants per group).
    ``Present'' in this table is counted at the group-task level (i.e., at least one
    valid record exists for that stream in that task).}
  \label{tab:coverage}
  \small
  \begin{tabular}{lccrrr}
    \toprule
    Data Type & Task & Expected & Present & Entries & Coverage \\
    \midrule
    \multirow{4}{*}{Task Outcomes} & T1 & 10 & 10 & 54 & 100\% \\
    & T2 & 10 & 10 & 94 & 100\% \\
    & T3 & 10 & 10 & 565 & 100\% \\
    & T4 & 10 & 10 & 272 & 100\% \\
    & \multicolumn{4}{l}{Total: 40/40 group–task pairs} & 100\% \\[3pt]
    \multirow{4}{*}{During-Task VAD} & T1 & 10 & 9 & 1146 & 90\% \\
    & T2 & 10 & 10 & 924 & 100\% \\
    & T3 & 10 & 10 & 832 & 100\% \\
    & T4 & 10 & 10 & 406 & 100\% \\
    & \multicolumn{4}{l}{Total: 39/40 group–task pairs (grp-09 T1 missing)} & 97.5\% \\[3pt]
    \multirow{4}{*}{Post-Task Surveys} & T1 & 10 & 10 & 2490 & 100\% \\
    & T2 & 10 & 10 & 2272 & 100\% \\
    & T3 & 10 & 10 & 2339 & 100\% \\
    & T4 & 10 & 10 & 2004 & 100\% \\
    & \multicolumn{4}{l}{Total: 40/40 group–task pairs} & 100\% \\[3pt]
    \multicolumn{5}{l}{\textbf{Grand total: 12,398 stimuli records}} \\
    \bottomrule
  \end{tabular}
\end{table}

\paragraph{How to read \Cref{tab:coverage}.}
For in-task probes and post-task questionnaires, completeness is near-perfect:
in-task VAD probes are 39/40 group-task pairs (97.5\%; one missing stream in grp-09 T1),
and post-task questionnaires are 40/40 (100\%).

\paragraph{Task Outcomes.}
All ten groups completed all four tasks and submitted formal decisions.
The elevated entry counts (e.g., 565 for T3 idea generation) reflect task structure:
multiple decision items per task (e.g., final\_format, final\_topic, selected\_candidate)
and per-participant contributions (e.g., individual token allocations in T4).

\paragraph{During-Task VAD.}
Valence, arousal, and dominance ratings were collected at multiple prompts
during each task. Coverage is near-complete (97.5\%); grp-09 experienced a technical
recording gap affecting T1 VAD data only. The high per-group entry counts
(83–127 on average) reflect the protocol's granular multi-prompt design,
providing temporal resolution into affect dynamics throughout each task.

\paragraph{Post-Task Surveys.}
Post-task survey protocols achieved 100\% group-level completion across all tasks.
Item-level completion rate exceeded 99\%; only two participants (grp-15 P1, grp-16 P2)
skipped optional demographic questions. Per-group entry counts (200–250) reflect
the full survey battery of task-dependent items (\textit{see} \cref{tab:survey-items}).

\paragraph{Missing Data Characterization.}
Item-level missing data rate is $< 1.3\%$ and exhibits random patterns
with no systematic bias by group, task, or participant position.
No session is entirely absent from any data type, and no group-task pair is
missing all records. The dataset does not require imputation for standard analyses.

%%=============================================================================
\section{Quality Control Audit}
\label{sec:qc-audit}

\begin{figure}[htb]
  \centering
  \includegraphics[width=\linewidth]{section5_missing_data_audit.png}
  \caption{Missing-data audit: proportion of expected records present per
    item and group. Darker cells indicate higher missing rates.}
  \label{fig:missing}
\end{figure}

\begin{figure}[htb]
  \centering
  \includegraphics[width=0.85\linewidth]{qc_audit_dashboard.png}
  \caption{QC audit dashboard covering data completeness, format consistency,
    deduplication, and synchronisation checks. Green\,=\,pass.}
  \label{fig:qc-dashboard}
\end{figure}

\paragraph{Interpretation.}
The stimuli and self-report data are largely complete across all ten groups.
Missing entries are item-level (individual non-responses) rather than
session-level gaps — no group is entirely absent from any core response
category.
The QC audit confirms format consistency and successful deduplication across
all groups, supporting the use of the full $N=10$ corpus for all analyses
reported in this document.

%%=============================================================================
\appendix

\section{grp-07 T1 Protocol Violation}
\label{app:grp07-protocol}

During task T1 (hidden-profile candidate selection), group-07 deviated from the 
intended protocol: two participants (P1 and P4) submitted the decision form instead 
of one designated representative.

\textbf{Findings:}
\begin{itemize}
  \item P1 submitted first at LSL~timestamp 328951.78 selecting \textit{Candidate C}
  \item P4 submitted ~2.7 seconds later, also selecting \textit{Candidate C}
  \item Both participants independently reached the same decision
\end{itemize}

\textbf{Impact on analysis:}
The duplicate submission does not affect the consensus outcome: all 10 groups 
(including grp-07) unanimously selected Candidate~C. Thus, T1 consensus remains 
$100\%$ (10/10 groups). The raw stimuli dataset retains both entries as recorded; 
researchers adhering to strict protocol compliance may exclude one submission 
(typically keeping P1's earlier submission).

\textbf{Data location:} 
Raw file at \texttt{sub-01\_ses-20260312\_grp-07\_run01\_task-T0T1T2T3T4\_stimuli\_answers.tsv}, 
rows corresponding to \texttt{task=T1}, \texttt{response\_type=form}, 
\texttt{item\_key=selected\_candidate}.

%%=============================================================================
\end{document}
